\documentclass[11pt]{article}

\usepackage[margin=1in]{geometry}
\usepackage{amsmath,amssymb,booktabs,hyperref}

\title{Cosmological Singularities from Entangled CFTs:\\
Resolution, Encoding, and Access}
\author{Stefano Antonini}
\date{April 9, 2026}

\begin{document}
\maketitle

\begin{abstract}
I revisit the current holographic literature on cosmologies generated by entangled conformal field theory states with one narrow question in mind: what, exactly, is established about the cosmological singularity? A source-level audit of five anchor papers shows a consistent pattern. The literature supports boundary sensitivity to interior cosmological dynamics, explicit Euclidean-to-Lorentzian constructions of cosmology and wormholes, and increasingly sharp statements about holographic encoding and final-state data. It does not yet supply a common boundary observable that rules out an unresolved big-bang or big-crunch singularity across the available model classes. The cleanest current reading is therefore conservative: the literature supports encoding and access claims more strongly than direct singularity-resolution claims.
\end{abstract}

\section{Introduction}

The recent literature on holographic cosmology has made it increasingly plausible that entangled CFT states can encode bulk regions with cosmological interpretation, sometimes behind horizons, sometimes on end-of-the-world or braneworld sectors, and in the latest work even for closed universes without asymptotically AdS boundaries \cite{Cooper:2018cmb,VanRaamsdonk:2021qgv,Antonini:2022ptt,Antonini:2024magnetic,Rath:2025baby}. The main conceptual danger is to slide too quickly from ``the CFT contains cosmological information'' to ``the singularity is resolved.'' Those are not the same claim.

The purpose of this note is therefore not to advertise a new singularity-resolution mechanism. It is to restate, as carefully as possible, what the verified anchor papers actually establish and what they still leave open. After checking the citation metadata against arXiv and published records, I keep only five references whose roles have already been fixed by the project artifacts. The resulting note is deliberately conservative and should be read as a corrected comparison paper rather than as a derivational advance.

\section{Verified Anchor Set}

\paragraph{Black hole microstate cosmology.}
The earliest anchor is the microstate construction of \cite{Cooper:2018cmb}. Its key result for the present discussion is operational rather than terminological: sufficiently large boundary subregions have entanglement entropy $S_A(t)$ that is sensitive to behind-the-horizon FRW evolution. This is genuine evidence that a boundary observable can probe interior cosmological dynamics. It is not, by itself, evidence that the singularity is absent.

\paragraph{Cosmology from confinement.}
Reference \cite{VanRaamsdonk:2021qgv} supplies the conceptual Euclidean-to-Lorentzian bridge. The point relevant here is structural: confinement-inspired Euclidean saddles can continue to cosmological or wormhole geometries. This paper is indispensable background, but it is not the sharpest source for a singularity-specific boundary diagnostic.

\paragraph{Accelerating cosmology from a holographic wormhole.}
Reference \cite{Antonini:2022ptt} sharpens the link between entangled boundary sectors and cosmological behavior. The paper relates cosmological observables to observables in the wormhole spacetime and shows how analytic continuation of the holographic wormhole data can produce an accelerating cosmology. The correction I need to keep explicit is that acceleration, or the mere existence of a cosmological continuation, is not yet a singularity criterion.

\paragraph{Magnetic braneworlds.}
The strongest open-universe microscopic anchor in the present set is \cite{Antonini:2024magnetic}. It constructs flat big-bang or big-crunch braneworld cosmologies and planar traversable wormholes as different analytic continuations of the same Euclidean saddle, and it gives an explicit BCFT strip description of the setup. This is the clearest microscopic realization in the current corpus, but even here the boundary data have not yet been organized into a universal observable that excludes an unresolved crunch.

\paragraph{Closed universes and holographic encoding.}
The most explicit encoding statement comes from \cite{Rath:2025baby}. With enough bulk entanglement, the closed-universe Hilbert space can be holographically encoded in the CFT, while in the no-entanglement limit direct holographic encoding breaks down. The paper also gives a dictionary between final-state data and CFT data, and shows how coarse-grained observables survive in the conventional effective description. The correction here is equally important: strong encoding and final-state control do not automatically imply direct microscopic reconstruction of the singular region.

\section{A Conservative Taxonomy}

To prevent repeated overclaiming, I organize the literature into three claim levels:
\begin{align}
\text{direct resolution} &: \exists\, O_{\mathrm{CFT}} \ \text{s.t.}\ O_{\mathrm{CFT}} \Rightarrow \text{no unresolved crunch}, \label{eq:resolution}\\
\text{encoded singularity} &: \mathcal{H}_{\mathrm{closed}} \hookrightarrow \mathcal{H}_{\mathrm{CFT}}, \label{eq:encoding}\\
\text{observer-limited access} &: \bar O_{\mathrm{CFT}} \ \text{survives while direct reconstruction fails}. \label{eq:access}
\end{align}
These definitions are intentionally schematic. Equation \eqref{eq:resolution} is a logical criterion, not a dynamical theorem; equation \eqref{eq:encoding} states an abstract encoding relation; equation \eqref{eq:access} isolates the weaker notion of coarse-grained or observer-dependent access.

This taxonomy also has the correct limiting behavior for the source material. In the large-subsystem regime of \cite{Cooper:2018cmb}, entanglement entropy becomes sensitive to interior FRW evolution, strengthening the case for observable access. In the zero-entanglement limit emphasized in \cite{Rath:2025baby}, direct holographic encoding fails, so a claim can at most remain in the observer-limited or final-state category. And in the Euclidean-to-Lorentzian continuations of \cite{VanRaamsdonk:2021qgv,Antonini:2024magnetic}, the existence of multiple continuations explains why the geometry alone cannot be promoted to a singularity-resolution verdict.

\section{Comparative Audit}

\begin{table}[t]
\centering
\small
\setlength{\tabcolsep}{4pt}
\begin{tabular}{p{2.3cm}p{3.0cm}p{4.0cm}p{4.0cm}}
\toprule
Reference & Observable basis & Strongest justified claim & Stronger reading not yet justified \\
\midrule
\cite{Cooper:2018cmb} & Large-subsystem entanglement entropy & Boundary sensitivity to behind-the-horizon FRW dynamics & Generic resolution of the FRW singularity \\
\cite{VanRaamsdonk:2021qgv} & Euclidean construction and Lorentzian continuation & Cosmology-wormhole-confinement bridge & Observable-level exclusion of an unresolved crunch \\
\cite{Antonini:2022ptt} & Cosmological observables related to wormhole observables & Entanglement-supported accelerating cosmology from a holographic wormhole & Acceleration as a singularity diagnostic \\
\cite{Antonini:2024magnetic} & BCFT strip path integral and braneworld continuation & Explicit microscopic realization of big-bang or big-crunch cosmology and wormhole from one Euclidean saddle & A universal crunch-resolving boundary observable \\
\cite{Rath:2025baby} & Encoding map, final-state dictionary, coarse-grained observables & Closed-universe encoding and final-state control with entanglement-dependent access & Direct microscopic reconstruction of the singular region in every regime \\
\bottomrule
\end{tabular}
\caption{Corrected comparison table after the citation and claim audit.}
\label{tab:audit}
\end{table}

Table \ref{tab:audit} makes the main correction to the earlier narrative. The literature does not support a uniform statement that entangled CFT cosmologies resolve their singularities. What it does support is a layered picture:
\begin{itemize}
\item one class of papers gives direct boundary sensitivity to interior time dependence,
\item another class gives explicit cosmology or wormhole constructions from the same Euclidean data,
\item the most recent closed-universe work gives the sharpest encoding statement and the clearest account of what survives when direct access fails.
\end{itemize}

Across those categories, the strongest robust conclusion remains comparative rather than absolute: encoding and access claims are presently stronger than direct resolution claims. This is exactly why the project-level synthesis result, \texttt{R-01-DIAGNOSTIC-SPLIT}, should remain provisional rather than being upgraded into a verified theorem.

\section{What Still Needs To Be Shown}

The missing ingredient is not another slogan about holography. It is a cross-model observable that does all three jobs in equation \eqref{eq:resolution}: it must see the crunch region, it must exclude an unresolved singularity rather than merely describe it, and it must survive the distinction between microscopic encoding and observer-limited access. None of the five anchor papers currently delivers that package across microstate, wormhole, braneworld, and closed-universe constructions at once.

That negative conclusion is scientifically useful. It identifies the real target for any follow-up phase: not a broader review, but an observable-by-observable comparison that can test whether the same boundary quantity can discriminate resolved from encoded-but-unresolved singular behavior. Until such a criterion exists, the corrected statement is that the current literature gives a precise language for access and encoding, but not yet a general proof of singularity resolution.

\section{Conclusion}

After a direct citation check and a claim-by-claim rewrite, the paper becomes much simpler and much harder to overread. The verified anchor set supports three distinct notions: direct observable sensitivity to cosmological dynamics, microscopic encoding of closed-universe data, and coarse-grained or final-state access when direct reconstruction is unavailable. Present evidence favors the second and third categories more strongly than the first. The corrected conclusion is therefore not that entangled CFTs have already resolved cosmological singularities, but that they have given us a sharper framework for distinguishing resolution from encoding and from mere observer access.

\bibliographystyle{unsrt}
\bibliography{references}

\end{document}
